\section*{अध्याय \ref{ch:b3:sensory-systems} --- संवेदी तंत्र}

\begin{solution}{exo:b3:sensory-systems:1}
तार से, संदेश से नहीं: हर संवेदी तंत्रिका कोई \omterm{def:b3:sensory-systems:principles}{नामांकित तार} है जो
मस्तिष्क के अपने ही क्षेत्र में समाप्त होती है, और दृक् तंत्रिका पर जो
भी पहुँचे उसे प्रकाश की तरह पढ़ा जाता है --- इसीलिए नेत्रगोलक पर दाब
कोई कौंध पैदा करता है और कान पर चोट कोई घंटी।
\end{solution}

\begin{solution}{exo:b3:sensory-systems:2}
अँधेरे में चक्रीय GMP धनायन वाहिकाएँ खुली रखता है और भीतर जाती कोई
स्थिर ``अँधेरे की धारा'' उस कोशिका को विध्रुवित तथा ग्लूटामेट छोड़ता
बनाए रखती है। प्रकाश \omterm{def:b3:sensory-systems:retina}{रोडॉप्सिन}, ट्रांसड्यूसिन और फ़ॉस्फ़ोडाइएस्टरेज़
सक्रिय करता है, जो cGMP नष्ट कर देती है; वाहिकाएँ बंद हो जाती हैं,
भीतर जाती धारा रुक जाती है, और वह कोशिका अतिध्रुवित होकर कम छोड़ती है
--- प्रकाश का संकेत किसी कमी से मिलता है।
\end{solution}

\begin{solution}{exo:b3:sensory-systems:3}
$I = I_{0}10^{L/10}$: \qty{20}{dB}, \qty{1e-10}{W/m^2}; \qty{60}{dB},
\qty{1e-6}{W/m^2}; \qty{110}{dB}, \qty{0.1}{W/m^2}। सबसे तेज़ से सबसे
धीमी: $10^{9}$।
\end{solution}

\begin{solution}{exo:b3:sensory-systems:4}
मीठा, उमामी और कड़वा G-प्रोटीन-युग्मित ग्राहियों से (एक मीठा, एक
उमामी, कोई पच्चीस कड़वे); नमकीन सोडियम वाहिका ENaC से; खट्टा किसी
प्रोटॉन वाहिका से। ``\omterm{def:b3:sensory-systems:chemical}{स्वाद}'' का अधिकांश उन वाष्पशील अणुओं की गंध है जो
मुँह से नाक तक पहुँचते हैं; नाक बंद हो तो केवल वे पाँच \omterm{def:b3:sensory-systems:chemical}{स्वाद} बचते हैं
और भोजन फीका लगता है।
\end{solution}

\begin{solution}{exo:b3:sensory-systems:5}
फ़ेख़नर: $S = (1/0.02)\ln(10\,000/100) = 50\ln 100 \approx 230$ पग।
गिनती से: $1.02^{n} = 100$ देता है $n = \ln 100/\ln 1.02 = 233$।
\end{solution}

\begin{solution}{exo:b3:sensory-systems:6}
$P = 1 - e^{-a}\sum_{k=0}^{5}a^{k}/k!$: $a = 2$: $0.017$; $a = 6$:
$0.55$; $a = 12$: $0.98$। आधी कौंधें $a \approx 5.7$ पर दिखती हैं;
\qty{10}{\%} अवशोषण के साथ उस कौंध को कॉर्निया पर लगभग $57$
फ़ोटॉन पहुँचाने पड़ेंगे।
\end{solution}

\begin{solution}{exo:b3:sensory-systems:7}
$55/3\times 1.3 \approx 24$; दाब के \omterm{def:b3:sensory-systems:cochlea}{डेसिबलों} में, $20\log_{10}24
\approx \qty{28}{dB}$। जुड़ी हुई हड्डियाँ वह कंपन संचरित नहीं कर
सकतीं, इसलिए ध्वनि \omterm{def:b3:sensory-systems:cochlea}{कर्णावर्त} तक केवल खोपड़ी से होकर चालन द्वारा
पहुँचती है, मध्य कान के लाभ के बिना: यानी कोई
\qtyrange{30}{50}{dB} की चालनी हानि।
\end{solution}

\begin{solution}{exo:b3:sensory-systems:8}
\omterm{def:b3:sensory-systems:cochlea}{रोम कोशिका} की वाहिका को \omterm{def:b3:sensory-systems:cochlea}{शीर्ष-कड़ी} सीधे खींचकर खोलती है --- कोई
संदेशवाहक नहीं, कोई एंज़ाइम नहीं --- इसलिए वह माइक्रोसेकंडों में
अनुक्रिया देती है; प्रकाश-ग्राही एक फ़ोटॉन को दो एंज़ाइमी चरणों से
प्रवर्धित करता है, जिनमें दसियों मिलीसेकंड लगते हैं। श्रवण को वह
समय-परिशुद्धि मिलती है जो अंतःकर्णी विलंबों से ध्वनियाँ ढूँढ़ने और कला
का पीछा करने को चाहिए; दृष्टि को एकल फ़ोटॉन गिनने की संवेदनशीलता मिलती
है, और वह गति में कीमत चुकाती है।
\end{solution}

\begin{solution}{exo:b3:sensory-systems:9}
$H(0.05) = -0.05\ln 0.05 - 0.95\ln 0.95 = 0.150 + 0.049 = 0.199$।
मनुष्य: $\ln\binom{400}{20} \approx 400\times 0.199 = 80$, यानी लगभग
$10^{34}$ प्रतिरूप। चूहा: $\ln\binom{1100}{55} \approx 1100\times 0.199
= 219$, यानी लगभग $10^{95}$।
\end{solution}

\begin{solution}{exo:b3:sensory-systems:10}
माध्य $500\times 0.1/40 = 1.25$ स्वतःस्फूर्त घटनाएँ प्रति खिड़की। $P(k
\le 5) = e^{-1.25}(1 + 1.25 + 0.78 + 0.33 + 0.10 + 0.03) = 0.998$,
इसलिए $P(k \ge 6) \approx 2\times 10^{-3}$। एक की देहली के साथ कोई
स्वतःस्फूर्त घटना अधिकांश खिड़कियों में होती ($1 - e^{-1.25} = 0.71$)
और अँधेरा कौंधों से भरा रहता; छह पर, झूठी चेतावनियाँ पाँच सौ खिड़कियों
में एक बार आती हैं। वह देहली वहाँ रखी है जहाँ \omterm{def:b3:sensory-systems:retina}{शलाकाओं} की संवेदनशीलता
नहीं, उनका शोर तय करता है।
\end{solution}

\begin{solution}{exo:b3:sensory-systems:11}
सबसे बड़ा विलंब $0.2/340 = \qty{590}{\micro s}$ (यानी किसी एक ओर से
आती ध्वनि)। मध्य-रेखा के पास $\Delta t \approx (d/c)\sin\theta$,
इसलिए $\qty{10}{\micro
s}$ का अर्थ है $\sin\theta = 0.017$, यानी लगभग \qty{1}{\degree}।
\qty{4}{kHz} से ऊपर आवर्तकाल (\qty{250}{\micro s}) उससे छोटा है जिससे
न्यूरॉन बँध सकें, और कुछ सौ माइक्रोसेकंड का विलंब एक से अधिक चक्र भर
जाता है, जिससे कला अस्पष्ट हो जाती है; तब मस्तिष्क दोनों कानों के बीच
के स्तर-अंतर और स्थान-कूट को काम में लेता है।
\end{solution}

\begin{solution}{exo:b3:sensory-systems:12}
किनारे से दूर, $r = I(1 - 2w)$: गहरी ओर $0.6$, चमकीली ओर $1.2$।
आख़िरी गहरे ग्राही पर ($I = 1$, पड़ोसी $1$ और $2$): $r
= 1 - 0.2\times 3 = 0.4$; पहले चमकीले पर: $r = 2 - 0.6 = 1.4$।
वह पग किसी न्यूनोच्छाल और किसी अत्युच्छाल में बढ़ा-चढ़ाकर दिखता है ---
यानी माख पट्टियाँ। \omterm{def:b3:sensory-systems:retina}{दृष्टिपटल} को विपर्यास तथा किनारे मिलते हैं, और
संचरित करने को छोटी गतिक परास; वह निरपेक्ष स्तरों के प्रति निष्ठा खोता
है, और चमक के भ्रम पैदा करता है।
\end{solution}

\begin{solution}{pb:b3:sensory-systems:1}
\textbf{1.} $0.01/3.6\times 10^{-19} = 2.8\times 10^{16}$ फ़ोटॉन प्रति
सेकंड।
\textbf{2.} पुतली का क्षेत्रफल $\pi(3.5\times 10^{-3})^{2} = 3.8\times 10^{-5}$
m$^{2}$। \qty{1}{km} पर: अंश $3.8\times 10^{-5}/1.26\times 10^{7} =
3\times 10^{-12}$, यानी प्रति सेकंड $8.5\times 10^{4}$ फ़ोटॉन।
\qty{10}{km} पर: प्रति सेकंड $850$।
\textbf{3.} \qty{0.1}{s} में \qty{10}{\%} अवशोषित होने पर:
\qty{1}{km} पर $850$, \qty{10}{km} पर $8.5$ --- दोनों छह से ऊपर; वह
मोमबत्ती दस किलोमीटर पर दिखती है (अधिकांश बार)।
\textbf{4.} प्रति खिड़की छह अवशोषित के लिए पुतली में प्रति सेकंड $600$
फ़ोटॉन चाहिए: $4\pi d^{2} = 3.8\times 10^{-5}\times 2.8\times 10^{16}/600
= 1.8\times 10^{9}$ m$^{2}$, यानी $d \approx \qty{12}{km}$। वहाँ $P(k \ge
6 \mid a = 6) = 0.55$: वह मोमबत्ती लगभग आधी खिड़कियों में दिखती है।
\textbf{5.} $e^{-6} = 0.0025$। देहली पर वह गिनती खिड़की-दर-खिड़की
उतार-चढ़ाव करती है --- $6 \pm 2.4$ --- इसलिए वह स्रोत आता-जाता लगता है:
यानी किसी मंद तारे की टिमटिमाहट (जिसमें वायुमंडल अपनी भी जोड़ देता है)।
\textbf{6.} पृष्ठभूमि का उतार-चढ़ाव, प्रति खिड़की $\sqrt{10^{4}} = 100$
फ़ोटॉन, मोमबत्ती के $8.5$ को डुबो देता है: संसूचन के लिए संकेत को
अँधेरी आँख की देहली नहीं, पृष्ठभूमि का शोर पार करना पड़ता है।
\textbf{7.} प्रति वाहिका $\qty{1}{pA}/200 = \qty{5}{fA}$; $5\times
10^{-15}/1.6\times 10^{-19} = 3\times 10^{4}$ आयन प्रति सेकंड।
\textbf{8.} तीसवाँ भाग; लगभग तीस एक साथ आते फ़ोटॉन हर वाहिका बंद कर
देंगे और उस \omterm{def:b3:sensory-systems:retina}{शलाका} को संतृप्त कर देंगे।
\textbf{9.} $10^{-12}\times 0.2 = 2\times 10^{-13}$ C, यानी $1.3\times 10^{6}$
धनायन: एक फ़ोटॉन दस लाख से अधिक आवेशों पर काबू रखता है।
\textbf{10.} वह \omterm{def:b3:sensory-systems:retina}{शलाका} मिलीसेकंडों में संतृप्त हो जाती है, हर वाहिका बंद,
और उसका \omterm{def:b3:sensory-systems:retina}{रोडॉप्सिन} बनने से तेज़ विरंजित होता है; \omterm{def:b3:sensory-systems:retina}{शंकु}, जिनका लाभ कम और
बंद होना तेज़ है, अनुक्रिया देते रहते हैं और अपनी परास प्रकाश के छह
दशकों भर खिसका लेते हैं।
\textbf{11.} ऊँचे लाभ को लंबा समाकलन चाहिए (हर चरण में समय लगता है और
अनुक्रिया जमा होने के लिए खिंचती है), इसलिए \omterm{def:b3:sensory-systems:retina}{शलाका} धीमी और संवेदी है;
\omterm{def:b3:sensory-systems:retina}{शंकु} का छोटा प्रवर्धन और जल्दी बंद होना \qty{20}{ms} में अनुक्रिया देते
हैं, जो गति के लिए काफ़ी तेज़ है, और कीमत यह कि उसे सौ फ़ोटॉन चाहिए।
\textbf{12.} $10^{8}/40 = 2.5\times 10^{6}$ स्वतःस्फूर्त समावयवन प्रति
सेकंड पूरे \omterm{def:b3:sensory-systems:retina}{दृष्टिपटल} भर --- पर प्रति खिड़की केवल $1.25$, यानी
$500$ \omterm{def:b3:sensory-systems:retina}{शलाकाओं} के किसी भी टुकड़े में, जो छह की देहली से कहीं नीचे है, इसलिए वह शोर
टुकड़े-दर-टुकड़े ठुकरा दिया जाता है; वह देहली और वह पूलन ठीक इसी के
लिए हैं।
\textbf{13.} \qty{0.1}{W/m^2}; \qty{5.5e-5}{m^2} पर, \qty{5.5e-6}{W}।
\textbf{14.} $5.5\times 10^{-6}\times 7200 = \qty{0.04}{J}$ --- यानी दो
घंटों में लगभग चार सौ वर्षा-बूँदों जितना।
\textbf{15.} $10^{-12}\times 5.5\times 10^{-5} = 5.5\times 10^{-17}$ W;
\qty{10}{ms} में, $5.5\times 10^{-19}$ J --- यानी लगभग एक या दो दृश्य
फ़ोटॉनों की ऊर्जा।
\textbf{16.} $0.3/5000 = 6\times 10^{-5}$ रेड, $0.003^{\circ}$; वह
विक्षेप लगभग तीन हाइड्रोजन परमाणुओं जितना है।
\textbf{17.} $10^{(110 - 85)/10} = 316$: \qty{110}{dB} पर दो घंटे
\qty{85}{dB} पर $630$ घंटों जितनी ऊर्जा पहुँचाते हैं --- यानी लगभग
महीने भर के कार्य-दिवस।
\textbf{18.} लगभग $120$ बस-ध्यानयोग्य पग, प्रति \omterm{def:b3:sensory-systems:cochlea}{डेसिबल} एक।
\textbf{19.} $2^{400} = 10^{400\log_{10}2} = 10^{120}$।
\textbf{20.} $\ln\binom{400}{20} \approx 400\,H(0.05) \approx 80$, जिसमें
से लगभग $2$ का स्टर्लिंग सुधार घटाकर: कोई $10^{34}$ प्रतिरूप।
\textbf{21.} वे ग्राही गिनिए जो एक प्रतिरूप में हों और दूसरे में नहीं: $5 +
5 = 10$, कुल $40$ में से (यानी कोई हैमिंग दूरी); जो प्रतिरूप तीन-चौथाई तक अतिव्यापी हों
वे लगभग वही \omterm{def:b3:sensory-systems:chemical}{ग्लोमेरुलस} चलाते हैं और लगभग वही गंध पढ़े जाते हैं।
\textbf{22.} प्रतिरूपों का $10^{12}/10^{34} = 10^{-22}$। सीमा वह कूट
नहीं है: ग्राही शोरभरे हैं, अनेक एक-सी सुर में बँधे हैं इसलिए उनकी
क्रियाशीलताएँ सहसंबद्ध हैं, असली अणु प्रतिरूपों का कोई छोटा उपसमुच्चय
ही बनाते हैं, और प्रतिरूपों का मस्तिष्कीय विभेदन तथा स्मरण सीमित है।
\textbf{23.} $\ln\binom{1100}{20} \approx 1100\,H(0.018) = 1100\times
0.091 \approx 100$, यानी लगभग $10^{43}$ --- यानी बीस-ग्राही प्रतिरूपों
की मानव संख्या से कोई $10^{9}$ गुना (और यदि वह चूहा प्रति गंध अधिक
ग्राही काम में ले तो कहीं अधिक)।
\textbf{24.} सामान्य गंध शायद ही बदलती है --- वह कूट अनावश्यक-बहुल है
और दूसरे ग्राही हर गंधद्रव्य ढो लेते हैं --- पर जो गंधद्रव्य कम
सांद्रता पर उसी ग्राही पर टिका हो वह या तो पकड़ में आना बंद कर देता है
या अपना चरित्र बदल देता है: यानी कोई विशिष्ट अगंधता, जैसी उन अनेक लोगों
में है जो ऐंड्रोस्टेनोन नहीं सूँघ सकते।
\textbf{25.} वह मोमबत्ती लगभग \qty{12}{km} तक दिखती है; देहली की ध्वनि
\qty{10}{ms} में कर्णपटह तक $5.5\times 10^{-19}$ J पहुँचाती है, यानी
एक-दो फ़ोटॉन; और कोई मानव नाक बीस-ग्राही प्रतिरूपों में से कोई
$10^{34}$ बना सकती है।
\end{solution}
